% !TEX root = ../main.tex
Federated Natural Policy Gradient (FedNPG) inherits the fast convergence of
Natural Policy Gradient (NPG) but
requires every agent to uplink its local Fisher information matrix $\bH_i \in
\Rb^{d\times d}$ each round. The resulting $\Oh(d^2)$ communication cost
becomes prohibitive for deep policies. \citet{lan2023fednpgadmm} reformulate
the global natural-gradient direction as a distributed quadratic program
solved by the Alternating Direction Method of Multipliers (ADMM), so that
each agent transmits only two $d$-dimensional vectors,
an $\Oh(d)$ cost, at the same convergence rate. We reproduce that method from
scratch in PyTorch and evaluate it on Swimmer-v4, Hopper-v4 and Humanoid-v4
with $N \in \{1,2,4,8\}$ agents over five random seeds. On Swimmer we match
the reported final rewards at two, four and eight agents and exceed them with
one; on Humanoid the method scales from
$993$ to $6329$ average return as agents increase from one to eight. On
Hopper, with a reduced configuration closer to the algorithm as written, four
and eight agents reach $2536$ and $2656$, within one standard deviation of the
reported values, and learn faster than one or two. Against the two baselines at eight agents, FedNPG-ADMM matches
standard-average FedNPG on Hopper and exceeds it on Humanoid. We further
quantify the communication
saving directly, measuring reductions by factors of $2.4\times 10^{3}$,
$9.2\times 10^{3}$ and $3.6\times 10^{5}$ in transmitted parameters per agent
per round on the three environments.
Finally, we quantize the two vectors each agent transmits, which the original
work does not attempt. Across five seeds, 8-bit stochastic quantization keeps
the final reward within the seed spread and the faster learning with more
agents on all three environments, for a further fourfold reduction in uplink against the
\texttt{float32} values exchanged (eightfold at the paper's 64-bit
accounting); at 6 bits Humanoid
collapses at every agent count.
